\documentclass[11pt]{article}

\usepackage[margin=1in]{geometry}
\usepackage{fancyhdr}

\title{\Large PLUS Session Outline \\ \normalsize \textit{MATH 1554:} Linear Algebra}
\author{Ahmed Mohammed}
\date{29th of August, 2023}

% ------- Set up header -------
\renewcommand{\sectionmark}[1]{\markright{#1}} % from https://tex.stackexchange.com/a/317393/227665
\pagestyle{fancy}
\fancyhead[L]{29th of August, 2023}
% \fancyhead[C]{\textit{\nouppercase{\rightmark}}}
\fancyhead[C]{\textit{MATH 1554 PLUS Session Outline}}
\fancyhead[R]{Mohammed \textbf{\thepage}}
% -------------------------------------------------------

\begin{document}
%\maketitle
\thispagestyle{fancy}
\section{\centering Open Activity}
	\begin{enumerate}
		\item \textit{Activity}. I will provide some questions on the slide that concerns definitions as a way to revise concepts covered in the past week and a half.
						I will provide the students some time to respond to the questions asked, and each group will elect a person to answer in behalf
						of that group.
		\item \textit{Collaboration}. Each member of a specific group will be discussing the definition amongst each other, so they'll have to work together
						to reach to an answer.
		\item \textit{Intended Goal}. The purpose of this activity is to serve as a revision for the concepts covered in previous lectures. Also, it'll help me
						understand where the students stand so that I may focus on some questions in the worksheet.
	\end{enumerate}

\section{\centering Main Activity}
	\begin{enumerate}
		\item \textit{Activity}. Students will be given a worksheet that goes over the previous 4 sections covered in lecture. Each group will work on an "activity",
						or rather a multi-part question. Once the allocated time is done, each group will elect a person to tell the whole session the answer
						they reached.
		\item \textit{Collaboration}. As students answer the questions, they'll have to work together. 
		\item \textit{Intended Goal}. The purpose of this activity is to test the students' understanding of concepts.
	\end{enumerate}

\section{\centering Closing Activity}
	\begin{enumerate}
		\item \textit{Activity}. In this activity, students will be asked two questions that foreshadow the next topic: homogenous and inhomogenuous equations,
			and how homogenous equations help us get closer to the solutions of the corresponding consistent inhomogenous system of linear equations.
		\item \textit{Collaboration}. \textsc{N/A}
		\item \textit{Intended Goal}. This activity will help students prepare for the next lecture and session.
	\end{enumerate}


\newpage
\fancyhead[L]{31st of August, 2023}
\section{\centering Open Activity}
	\begin{enumerate}
		\item \textit{Activity}. I will begin the session by providing a general review of the topics covered in lectures. The way I'll do this is by assigning groups and then
						asking each group to answer a question. A different group might have to \emph{justify} the group's answer as well. Groups will have
						time to discuss amongst each other.
		\item \textit{Collaboration}. During the allocated time the groups are allowed to discuss amongst each other, the students will have a chance to introduce each other
						quickly and answer the posed question.
		\item \textit{Intended Goal}. Before jumping into the worksheet, we will be able to revise topics covered in lecture.
	\end{enumerate}

\section{\centering Main Activity}
	\begin{enumerate}
		\item \textit{Activity}. Students will be given a worksheet that goes over the previous section covered in lecture. Each group will work on an "activity",
						or rather a multi-part question. Once the allocated time is done, each group will elect a person to tell the whole session the answer
						they reached. I'll then ask a group to reveal the answer they reached and to explain why that's chose to do so. If they couldn't, a different
						group will intervene and provide their own justification.
		\item \textit{Collaboration}. As students answer the questions, they'll have to work together. 
		\item \textit{Intended Goal}. The purpose of this activity is to test the students' understanding of concepts.
	\end{enumerate}

\section{\centering Closing Activity}
	\begin{enumerate}
		\item \textit{Activity}. We'll be foreshadowing the next main idea of the course: \textit{linear independence}.
		\item \textit{Collaboration}. \textsc{N/A}
		\item \textit{Intended Goal}. This activity will help students prepare for the next lecture and session.
	\end{enumerate}

\newpage
\fancyhead[L]{5th of September, 2023}
\section{\centering Open Activity}
	\begin{enumerate}
		\item \textit{Activity}. The session will begin with a general review of the topic of linear independence. Each group will have to answer a specific set of
			of questions, and then we will go over them after an allocated amount of time.
		\item \textit{Collaboration}. Students will be expected to go over the questions together, correcting each other as they go.
		\item \textit{Intended Goal.} This activity is meant to be a review of the foundation on which the rest of the session is built on.
	\end{enumerate}'

\section{\centering Main Activity}
	\begin{enumerate}
		\item \textit{Activity.} Students will answer 3 short answer questions and 5 more questions which go over the main topic of the session.
		\item \textit{Collaboration.} Students will be expected to work together.
		\item \textit{Intended Goal.} Students will be able to deepen their understanding of the topic covered in session. They'll also be able to ask questions 
			as well.
	\end{enumerate}
\section{\centering Closing Activity}
	\begin{enumerate}
		\item \textit{Activity.} We will hint into the next main idea which deals with linear transformations.
		\item \textit{Collaboration.} N/A
		\item \textit{Intended Goal.} Students will be able to revise for the next session, especially as the exam comes closer. 
	\end{enumerate}

\newpage
\fancyhead[L]{7th of September, 2023}
\section{\centering Open Activity}
	\begin{enumerate}
		\item \textit{Activity.} I will review the material with the students as they simultaneously form groups.
		\item \textit{Collaboration.} Students will begin to form groups and will disuss the review questions together.
		\item \textit{Intended Goal.} This is meant to refresh the concepts and ideas covered in previous lecture(s).
	\end{enumerate}'

\section{\centering Main Activity}
	\begin{enumerate}
		\item \textit{Activity.} Students will go over the created worksheet. After some allocated time, I'll review their answers
		\item \textit{Collaboration.} Students will be expected to go over the questions together during the allocated time.
		\item \textit{Intended Goal.} Students will be able to test their understanding of the contents covered in lectures.
	\end{enumerate}
\section{\centering Closing Activity}
	\begin{enumerate}
		\item \textit{Activity.} We will highlight the next main idea: standard matrices.
		\item \textit{Collaboration.} N/A
		\item \textit{Intended Goal.} This activity will give the students a glimpse as to what they should expect.
	\end{enumerate}

\newpage
\fancyhead[L]{19th of October, 2023}
\section{\centering Open Activity}
	\begin{enumerate}
		\item \textit{Activity.} I will review the material with the students as they simultaneously form groups.
		\item \textit{Collaboration.} Students will begin to form groups and will disuss the review questions together.
		\item \textit{Intended Goal.} This is meant to refresh the concepts and ideas covered in previous lecture(s).
	\end{enumerate}'

\section{\centering Main Activity}
	\begin{enumerate}
		\item \textit{Activity.} Students will go over the created worksheet. After some allocated time, I'll review their answers
		\item \textit{Collaboration.} Students will be expected to go over the questions together during the allocated time.
		\item \textit{Intended Goal.} Students will be able to test their understanding of the contents covered in lectures.
	\end{enumerate}
\section{\centering Closing Activity}
	\begin{enumerate}
		\item \textit{Activity.} We will go over the hardest problem on the set, and allow students to discuss it amongst themselves.
		\item \textit{Collaboration.} N/A
		\item \textit{Intended Goal.} This activity is meant to help students evaluate their own understanding of the topic covered.
	\end{enumerate}

\newpage
\fancyhead[L]{24th of October, 2023}
\section{\centering Open Activity}
	\begin{enumerate}
		\item \textit{Activity.} We will begin the session by reviewing key concepts and identifying what the weakest comprehension points are.
		\item \textit{Collaboration.} Students are expected to be engaged in the process and explain concepts as \emph{other} students ask questions.
		\item \textit{Intended Goal.} The intention of this activity is to serve as a way to allow me to assess the understanding of the students.
	\end{enumerate}

\section{\centering Main Activity}
	\begin{enumerate}
		\item \textit{Activity.} Students will be given a certain amount of time to answer questions on the worksheet. After that, we'll go over the answers 
			together.
		\item \textit{Collaboration.} Students will answer the questions as they discuss it with their group.
		\item \textit{Intended Goal.} Students will be given the opportunity to assess their own understanding of the material.
	\end{enumerate}
\section{\centering Closing Activity}
	\begin{enumerate}
		\item \textit{Activity.} We'll go over the hardest problem on the worksheet.
		\item \textit{Collaboration.} N/A
		\item \textit{Intended Goal.} The goal of this activity is to allow students to go over the hardest problem.
	\end{enumerate}

\newpage
\fancyhead[L]{26th of October, 2023}
\section{\centering Open Activity}
	\begin{enumerate}
		\item \textit{Activity.} We will begin the session by reviewing our answers for the worksheet from the previous session.
		\item \textit{Collaboration.} Students are expected to be engaged in the process and explain concepts as \emph{other} students ask questions.
		\item \textit{Intended Goal.} The intention of this activity is to allow us to continue off from where we picked off last session.
	\end{enumerate}

\section{\centering Main Activity}
	\begin{enumerate}
		\item \textit{Activity.} After revision, students will work through rest of the problems together. At the same time, I'll be answering questions regarding
			the previous questions on the worksheet to students that might be facing difficulty with the previous problems.
		\item \textit{Collaboration.} Students will answer the questions as they discuss it with their group.
		\item \textit{Intended Goal.} Students will be given the opportunity to assess their own understanding of the material.
	\end{enumerate}
\section{\centering Closing Activity}
	\begin{enumerate}
		\item \textit{Activity.} We'll go over the hardest problem on the worksheet.
		\item \textit{Collaboration.} N/A
		\item \textit{Intended Goal.} The goal of this activity is to allow students to go over the hardest problem.
	\end{enumerate}

\newpage
\fancyhead[L]{26th of October, 2023}
\section{\centering Open Activity}
	\begin{enumerate}
		\item \textit{Activity.} We will go through a desmos plot I did to get a better understanding of orthogonal projections.
		\item \textit{Collaboration.} N/A.
		\item \textit{Intended Goal.} The intention of this activity is to gives students a better understanding of orthogonal projection.
	\end{enumerate}

\section{\centering Main Activity}
	\begin{enumerate}
		\item \textit{Activity.} I'll provide students an allocated amount of time to finish each section of the worksheet.
		\item \textit{Collaboration.} Students will answer the questions as they discuss it with their group.
		\item \textit{Intended Goal.} Students will be given the opportunity to assess their own understanding of the material.
	\end{enumerate}
\section{\centering Closing Activity}
	\begin{enumerate}
		\item \textit{Activity.} We'll go over the hardest problem on the worksheet, which is the question before the last.
		\item \textit{Collaboration.} N/A
		\item \textit{Intended Goal.} The goal of this activity is to allow students to go over the hardest problem, or to ask questions.
	\end{enumerate}
\end{document}